\section{引言}\label{sec:intro}

虚假新闻检测的一个基本设定是只根据新闻文本本身判断真伪：没有传播数据，也不检索外部证据。这一\emph{闭卷}设定对应新闻刚出现、尚无事实核查报道与社交反馈的阶段。长期以来，这一任务由微调的小规模预训练语言模型（small language model, SLM）承担\citep{shu2017fake,devlin2019bert}。大语言模型（LLM）出现后，一个自然的想法是利用其常识与推理能力。\citet{hu2024bad} 发现，LLM 直接判断真伪不如微调的 BERT，但它写出的分析理由可以作为 SLM 的补充输入，由此形成了“LLM 提供理由、SLM 做裁决”的协同范式\citep{hu2024bad,he2026factguard,liu2024teller,nan2024genfend}。另一条路线让多个 LLM 智能体就新闻的真伪展开辩论\citep{liu2025ted,han2025d2d}。

这些方法有两个共同的薄弱环节。第一，LLM 的理由并不总是可用，现有方法对“何时采信 LLM”的建模停留在整条理由一级：ARG 与 FactGuard 为每条理由预测一个权重，监督信号是“这条理由对应的 LLM 判断是否等于标签”。一条理由内部哪些内容可信、哪些是幻觉，无从区分。第二，辩论类方法每条新闻需要多次 LLM 调用，而其消融实验显示贡献主要来自可训练的裁决模块，而非辩论文本本身\citep{liu2025ted}。

\paragraph{本文的方法}
针对这两点，本文提出 \method{}。LLM 在一次结构化调用中把新闻分解为若干原子声明，并为每条声明同时写出最强的支持论证与质疑论证，相当于把辩论压缩到声明一级。可训练的裁决器把一条新闻看作由这些声明构成的包，用多示例学习聚合为新闻级判定，训练只需新闻级标签。裁决器包含两个门控：\emph{相关性门控}判断一个分析单元是否真的属于这条新闻，不属于的单元成为不携带证据的空实例，全部为空时预测精确地退回到只依赖新闻文本的分支；\emph{判定门控}调节 LLM 证据的权重。相关性门控用训练时构造的错配负例监督。

\paragraph{本文的评测}
在评测上，本文力求排除几个常见的混淆因素：所有方法使用同一个开源 LLM（Qwen2.5-7B）重新生成输入；所有方法在同一组学习率上搜索；主实验报告 5 个随机种子的均值、标准差与显著性检验。本文还加入了两个在已有工作中缺席的基线：把同一个 LLM 用 LoRA 直接微调为分类器，以及在 BERT 的输出上简单加一项 LLM 零样本评分。在分布内精度之外，本文设计了四项受控测试：LLM 分析被整条或部分替换、保持事实的风格改写、位置已知的单句事实篡改，以及留一领域。

\paragraph{主要发现}
实验结果并不都支持本文最初的设想，我们如实报告。
\begin{enumerate}
\item \textbf{分布内精度持平。} \method{} 在 Weibo21 与 GossipCop 上的 macro-F1 为 \ptwo{CN/main/clara/macf1} 与 \ptwo{EN/main/clara/macf1}，与 ARG（\ptwo{CN/main/arg/macf1} 与 \ptwo{EN/main/arg/macf1}）没有显著差异。BERT 加零样本评分这一最简单的做法达到 \ptwo{CN/main/bert_zs/macf1} 与 \ptwo{EN/main/bert_zs/macf1}，不逊于任何一种让 SLM 阅读 LLM 理由的结构。在多种子设置下，已发表工作之间约 1 个点的差距落在随机波动之内。
\item \textbf{相关性门控有效，适用范围有限。} 它区分“属于本新闻的分析单元”与“被替换进来的单元”的 AUC 超过 0.98，使 \method{} 在分析被替换时退回到接近纯文本模型的水平；去掉它，在 Weibo21 上模型在同样条件下比 BERT 低 \pone{CN/derived/no_relevance_gap} 个点。在 Weibo21 上，当分析中的声明全部被逐条替换时，\method{} 比 ARG 高 \ptwo{CN/noise_sig/unit100/arg/diff_mean} 个点（对测试样本的检验显著，对随机种子的检验处于边缘）；GossipCop 上没有这一优势。ARG 本身对整条替换已相当稳健，给它加上同样的错配负例后，在整条替换下与 \method{} 相当。
\item \textbf{“LLM 是否可靠”在闭卷条件下学不出来。} 判定门控跟随的是新闻标签，而不是 LLM 判断的正误。ARG 的理由可用性预测器同样没有提供分类器自身预测之外的信息。
\item \textbf{所有闭卷检测器都主要依赖文风。} 把假新闻改写成通稿文风、真新闻改写成营销号文风后，Weibo21 上所有 SLM 类方法的 macro-F1 降到 51--53。LLM 自身也是如此：其零样本判断的 AUC 从 \pone{CN/style_llm/zero-shot P(fake)/auc_orig} 降到 \pone{CN/style_llm/zero-shot P(fake)/auc_style}。把分析拆到声明一级并不能改变这一点。
\item \textbf{对局部事实改动不敏感。} 把真新闻中的一句话改成与原文事实不符的陈述，各方法的预测几乎不变；声明级的定位只略高于随机。
\end{enumerate}

\paragraph{贡献}
（1）提出声明级双立场论证与带相关性门控的多示例裁决框架 \method{}，以及不调用 LLM 的学生模型 \student{}；（2）给出一套控制了 LLM、超参数搜索与随机性的评测流程，补充了两个被忽略的基线与四项受控测试；（3）报告了上述正面与负面的发现，其中“闭卷条件下 LLM 的分析与基于它的检测器都主要依赖文风”这一点，对该范式的后续研究有直接的警示意义。代码与全部实验配置已公开\footnote{\url{https://github.com/roman0905/paper}}。
